\section{Philosphy}
This user manual is part of a bigger project called Fort-Sovereign (edgy name, I know...).
The intent of this manual is to make anyone, whoever it is, capable of creating a server in their own home.
The objective is to allow even more people to get away from corporations that favor shareholders rather than their clients.
So, the responsibility falls to us, to show them that they are not the only option. 

The goal is that, by the end of this manual, you will be able to create and deploy your very own server. 
You will learn all the surface-level knowledge required to do so. \\


\subsection{License}
Both this user manual as well as the project repository is licensed under the \textbf{GPLv3} distribution.
Anyone is free to take, change and redistribute both this manual and the repository content under the same license \textit{only}.

\section{Architecture}
This section details the specifications of my home server. 
It lists all hardware, software and external resources. 
It also lists all services I intend to use to deploy all my services. 
For exemple, 


\subsection{Hardware}
The hardware is part of an old gaming PC, with only the HDDs added to it.
\subsubsection{CPU \& GPU}
The CPU is i5-7400 @ 3GHz. It is efficient and more than enough for the tasks at hand.

The GPU is a GTX 1060 6GB Nvidia graphics card. It is essential for transcoding, machine learning and other such tasks. It is also more than enough.

\subsubsection{Storage}
The server utilises three 8\unit{To} HDDs (SEAGATE BarraCuda) in a Raidz-1 pool configuration.
They host all of the files which are unnecessary to the running of the server OS.

I am also using an old 250\unit{Go} SSD, which is the boot drive of my server. 
It holds no data aside from the actual OS itself. Again, it is \textit{old} and I am not about to rely on it.

Additionally, I have another, much newer 250\unit{Go} SSD, which I am going to use to installing the new OS and then possibly as a cache for large file handling (the kind that Unmanic does, more on that later)

\subsection{Software}
This part lists the software I plan on using that are not services, but are necessary to the function of the server. It does not include external services.

\subsubsection{Operating System}
The operating system is \textbf{Debian 13 (Trixie)}\cite{debian_wiki}. 
It is a robust, stable and complete linux distro that is perfect for a simple server environment.
It is also possible to move from debian towards \textbf{Proxmox VE 9.4} by changing the source repository of the OS. If the server ever becomes part of cluster, this will facilitate the migration.

\subsubsection{Service Deployment}
The primary method of deploying services is using \textbf{Docker} containers to run the services.
This will be combined with \textbf{Docker Compose} to easily manage the services and stacks by using a git repository
The obvious choice is here is Docker, but combined with Docker Compose.
This setup allows me to configure the *.yaml files from a git repository. 
This method is also familiar to me and is ideal for my intended Git-Ops workflow.



\subsubsection{Storage Management}
\gls{zfs} is possibly the defining trait of the storage side of my server.
It is a very powerful service that can create pools (arrays of multiple HDDs) and convert them into a
R.A.I.D (Redundant Array of Independent Disks) configuration.


Specifically, my 'tank' storage pool is in a RAID-Z1 configuration.
It allows a single drive failure without losing any data. So if one drive dies, I do not lose all of the data on my server.
However, it is not a backup. That will be explored further below.

It also enables a bunch more things, but it will be explored later.

Despite starting from scratch, ZFS allows me to do one extremely powerful thing :
\begin{lstlisting}[language=python]
    zfs export tank
\end{lstlisting}

Using this command allows me to prepare my disks to be transfered to a new host. 
In my case, it does not change machines but since I will be re-installing the OS, it is a necessary safety measure.

\subsubsection{Dataset Hierarchy}
The \texttt{tank} pool is partitioned into logical datasets to apply specific ZFS properties optimized for the underlying data type.

\begin{table}[h]
    \centering
    \begin{tabular}{|l|l|l|l|}
        \hline
         \textbf{tank} & \textbf{Purpose} & \textbf{Recordsize} (128K) & \textbf{Snapshot Policy} \\ \hline
         tank/appdata & Configurations & 128K & Daily \\
         tank/db & Databases & 8K & Hourly \\
         tank/media & Photos / Videos / Music & 1M & Daily \\
         tank/backups & Local Backups & 128K & Hourly \\ \hline
    \end{tabular}
    \caption{ZFS Datasets}
    \label{tab:placeholder}
\end{table}

\subsection{Necessary Services}
This part details the services that a are necessary to run and deploy (almost) all the services I intend on using.
It will not list 'User-Oriented' services, but rather those necessary to their deployment. 
